% abstract.tex -- abstract draft for the arXiv v1 paper.
% Source: research/02-paper-outline.md §2 (v0.2, 2026-08-18), copied
% verbatim (LaTeX-escaped). ~150 words.
%
% TODO (W2): recheck every number against the final run before
% submission -- the 3,500/3,500 cell counts and the metric list must
% match the released results (research/01-experiment-protocol.md §2;
% 05-arxiv-readiness.md checklist F).

Text watermarking is the primary mechanism proposed for EU AI Act
Art.~50 transparency obligations on machine-generated content. We
measure its robustness against realistic user-side editing. Using
same-config detection over 3,500 watermarked and 3,500 unwatermarked
generations (KGW, SynthID-Text, EXP, Unigram, and SIR schemes, in
English and German/French/Spanish), we evaluate a layered removal
pipeline that combines formatting-layer cleanup (invisible Unicode,
bidi) with statistical rewriting driven by detection feedback, and we
report ROC-based metrics (AUROC, TPR@1\%FPR) and quality metrics
(perplexity, BERTScore, ROUGE-L) at the point of detection collapse. We
find that quality-preserving paraphrase and translation round-trip
collapse KGW-class detection to near chance, that SynthID-Text resists
token substitution but not paraphrase, that layering dominates single
attacks at equal quality cost, and that multilingual texts are
systematically more fragile. We discuss implications for Art.~50
compliance and release our corpus, configs, and harness.
